\documentclass[UTF8,zihao=5,twoside]{ctexart}
\usepackage[a4paper,top=2.5cm,bottom=2.5cm,left=3cm,right=3cm,headsep=1.5cm,footskip=1.5cm,bindingoffset=0cm]{geometry}
\usepackage{array}
\usepackage{tabularx}
\usepackage{longtable}
\usepackage{multirow}
\usepackage{makecell}
\usepackage{fancyhdr}
\usepackage{setspace}
\usepackage{ragged2e}
\usepackage{calc}

\setCJKmainfont[AutoFakeBold=2.5]{SimSun}
\setmainfont{Times New Roman}
\pagestyle{fancy}
\fancyhf{}
\renewcommand{\headrulewidth}{0pt}
\renewcommand{\footrulewidth}{0pt}
\fancyfoot[C]{\zihao{5}中国高校计算机大赛-人工智能创意赛（鸿蒙赛道）组委会编\hspace{2.0em}2026.06}
\fancyfoot[RO,LE]{\zihao{5}\thepage}

\setlength{\parindent}{2em}
\setlength{\parskip}{0pt}
\setstretch{1.0}
\renewcommand{\arraystretch}{1.35}
\setlength{\tabcolsep}{2.5pt}

\newcommand{\blankline}{\vspace{\baselineskip}\par}
\newcommand{\centertitle}[2]{{\centering\zihao{#1}\bfseries #2\par}}
\newcommand{\maintitle}[1]{{\centering\zihao{2}\bfseries #1\par}}
\newcommand{\authorline}[1]{{\centering\zihao{3}\bfseries #1\par}}
\newcommand{\firstheading}[1]{{\noindent\zihao{3}\bfseries #1\par}}
\newcommand{\secondheading}[1]{{\noindent\zihao{4}\bfseries #1\par}}
\newcommand{\thirdheading}[1]{{\noindent\zihao{-4}\bfseries #1\par}}
\newcommand{\fourthheading}[1]{{\noindent\zihao{5}\bfseries #1\par}}
\newcommand{\centerlinez}[2]{{\centering\zihao{#1}#2\par}}
\newcommand{\leftlinez}[2]{{\noindent\zihao{#1}#2\par}}
\newcommand{\rightlinez}[2]{{\raggedleft\zihao{#1}#2\par}}
\newcommand{\tblcell}[1]{\makecell[c]{#1}}
\newcommand{\tbllcell}[1]{\makecell[l]{#1}}

\begin{document}

\thispagestyle{fancy}
\maintitle{2026中国高校计算机大赛—人工智能创意赛}
\blankline
\maintitle{鸿蒙赛道作品说明文档（初赛）}

\vspace{3.7cm}

\centerlinez{4}{参赛学校：\underline{\makebox[7.8cm][c]{西安交通大学}}}
\centerlinez{4}{团队名称：\underline{\makebox[7.8cm][c]{枣La}}}
\centerlinez{4}{作品名称：\underline{\makebox[7.8cm][c]{《枣寻（LocateS）》}}}
\centerlinez{4}{赛题方向：\underline{\makebox[7.8cm][c]{应用创新}}}
\blankline
\centerlinez{4}{联系人（队长）：\underline{\makebox[6.5cm][c]{伍洲宜}}}
\centerlinez{4}{联系电话（队长）：\underline{\makebox[6.3cm][c]{18165717179}}}


% \vspace{4.9cm}

% \centertitle{3}{作品说明文档（初赛）}
% \centertitle{3}{包含但不限于以下内容：}
% \blankline

% \begin{enumerate}
%   \itemsep0pt
%   \item 参赛团队信息表
%   \item 作品原创性声明
%   \item 创意描述（30字以内）
%   \item 设计稿/技术方案
%   \item 介绍文档（800字以内）
% \end{enumerate}

% \vfill
\clearpage

\centertitle{3}{2026中国高校计算机大赛—人工智能创意赛}
\centertitle{3}{参赛团队信息表}

\begingroup
\scriptsize
\centering
\begin{longtable}{|>{\centering\arraybackslash}p{1.2cm}|*{12}{>{\centering\arraybackslash}p{0.86cm}|}}
\hline
作品名称 & \multicolumn{12}{c|}{《枣寻（LocateS）》} \\
\hline
团队名称 & \multicolumn{12}{c|}{枣La} \\
\hline
参赛学校 & \multicolumn{12}{c|}{西安交通大学} \\
\hline
\makecell{参赛\\赛题方向} & \multicolumn{12}{p{9.6cm}|}{（$\surd$）应用创新 （\,）Agent创新  （）用户体验创新 （）操作系统智能创新}\\ %\quad *请在所选赛题方向括号内打钩，示例：（$\surd$）应用创新} \\
\hline
\multicolumn{13}{|c|}{团队队员基本信息（可跨校、跨专业组队）} \\
\hline
姓名 & 学校全称 & \multicolumn{2}{c|}{院（系）全称} & \multicolumn{2}{c|}{专业全称} & 年级 & 毕业时间 & \multicolumn{2}{c|}{联系电话} & \multicolumn{2}{c|}{邮箱} & 团队分工 \\
\hline
伍洲宜 & 西安交通大学 & \multicolumn{2}{c|}{软件学院} & \multicolumn{2}{c|}{软件工程} & 大二 & 2028.6 & \multicolumn{2}{c|}{18165717179} & \multicolumn{2}{c|}{2805776721@qq.com} & 队长 \\ 
\hline
徐沈茁 & 西安交通大学 & \multicolumn{2}{c|}{机械工程学院} & \multicolumn{2}{c|}{智能制造工程} & 大一 & 2029.6 & \multicolumn{2}{c|}{13588493487} & \multicolumn{2}{c|}{xushenzhuo@stu.xjtu.edu.cn} & 队员 \\

\hline
\multicolumn{13}{|c|}{团队指导教师信息（指导老师须与队长同校）} \\
\hline
姓名 & \multicolumn{2}{c|}{院（系）全称} & \multicolumn{2}{c|}{职称} & \multicolumn{3}{c|}{研究方向} & \multicolumn{2}{c|}{联系电话} & \multicolumn{3}{c|}{联系邮箱} \\
\hline
陈** & \multicolumn{2}{c|}{信息科学技术学院} & \multicolumn{2}{c|}{教授} & \multicolumn{3}{c|}{深度学习、自然语言处理} & \multicolumn{2}{c|}{139****310} & \multicolumn{3}{c|}{32***@qq.com} \\
\hline
\multicolumn{13}{|c|}{团队成员优势描述} \\
\hline
\multicolumn{13}{|p{12.0cm}|}{\parbox[t][4.2em][t]{10.8cm}{可列举描述团队\\（1）成员个人或集体重要学术成果或项目经历；\\（2）各成员的擅长领域、分工和互补情况}} \\
\hline
\end{longtable}
\endgroup

% \vspace{0.2cm}

\newpage
\centertitle{3}{\underline{\hspace{1.5em}《枣寻（LocateS）》\hspace{1.5em}} 作品原创性声明}
\blankline

\leftlinez{5}{郑重声明：承诺本参赛队伍报名信息真实有效；呈交的参赛作品相关资料以及所完成的作品实物等相关成果，是本团队独立进行研究工作所取得的成果，除文中已经注明引用的内容外，本作品说明文档不包含任何其他个人或集体已经发表或撰写过的作品成果，不侵犯任何第三方的知识产权或其他权利。本声明的法律结果由本参赛队承担。}
\blankline\blankline
\leftlinez{5}{参赛队员签名（团队全部成员）：}
\blankline\blankline\blankline
\rightlinez{5}{日期：\hspace{2em}年\hspace{2em}月\hspace{2em}日}
\blankline\blankline
\leftlinez{5}{指导老师审核签名：}
\blankline\blankline\blankline
\rightlinez{5}{日期：\hspace{2em}年\hspace{2em}月\hspace{2em}日}
\blankline\blankline
\leftlinez{5}{（注：本页签名可以打印纸质文件后签名，提供扫描件；或者使用电子签名。不论哪种方式，需保证页面内容完整、字迹清晰，请与本项目创意书作为同一文档提交，否则项目创意书提交可能无效。）}

\clearpage

\maintitle{作品说明文档}
\blankline

\secondheading{创意描述}

空间影像标记定位物品，事件清单驱动智能收纳提醒

\secondheading{设计稿}

\blankline

\thirdheading{一、系统架构}

\vspace{0.5cm}

\begin{center}
\fbox{
\begin{tabular}{|>{\centering\arraybackslash}p{3.2cm}|>{\centering\arraybackslash}p{3.2cm}|>{\centering\arraybackslash}p{3.2cm}|>{\centering\arraybackslash}p{3.2cm}|}
\hline
\multicolumn{4}{|c|}{\zihao{5}\bfseries 枣寻（LocateS）系统架构} \\
\hline
\multicolumn{4}{|c|}{\zihao{5}\bfseries 页面层（Pages）} \\
\hline
首页（物品Tab） & 空间管理页 & 事件管理页 & 个人中心 \\
添加物品页 & 物品详情页 & 事件详情页 & 背景设置页 \\
\hline
\multicolumn{4}{|c|}{\zihao{5}\bfseries 组件层（Components）} \\
\hline
\multicolumn{2}{|c|}{物品卡片 / 分类网格} & \multicolumn{2}{c|}{事件卡片 / 空状态} \\
\hline
\multicolumn{2}{|c|}{确认删除对话框} & \multicolumn{2}{c|}{页面背景（渐变/图片）} \\
\hline
\multicolumn{4}{|c|}{\zihao{5}\bfseries 数据层（Model + Repository + Database）} \\
\hline
\multicolumn{2}{|c|}{Item / Space / Event / Category 模型} & \multicolumn{2}{c|}{ItemRepo / SpaceRepo / EventRepo} \\
\hline
\multicolumn{4}{|c|}{Database 单例 — RelationalStore（5张表）} \\
\hline
\multicolumn{4}{|c|}{\zihao{5}\bfseries 平台层} \\
\hline
\multicolumn{4}{|c|}{HarmonyOS 6.1 (API 24) · ArkTS · ArkUI · HDS · systemShare · 服务卡片} \\
\hline
\end{tabular}
}
\end{center}

\blankline
\thirdheading{二、核心交互流程}

\vspace{0.3cm}

\begin{center}
\fbox{
\begin{tabular}{|>{\centering\arraybackslash}p{2.8cm}|>{\centering\arraybackslash}p{5.5cm}|>{\centering\arraybackslash}p{5.0cm}|}
\hline
\zihao{5}\bfseries 模块 & \zihao{5}\bfseries 用户操作 & \zihao{5}\bfseries 系统响应 \\
\hline
\multirow{4}{*}{物品管理} & 搜索/分类筛选物品 & 实时过滤物品列表 \\
\cline{2-3}
& 点击"+"添加物品（半模态面板） & 写入item表，通知列表刷新 \\
\cline{2-3}
& 点击物品进入详情页 & 展示全部字段，可修改分类/位置 \\
\cline{2-3}
& 管理分类（添加/删除） & 动态更新category表，级联更新物品 \\
\hline
\multirow{5}{*}{空间管理} & 创建空间（上传平面图/拍照） & 写入space表（level=1），存储图片 \\
\cline{2-3}
& 点击图片标记分区位置 & 记录markerX/Y坐标，生成彩色标记点 \\
\cline{2-3}
& 拖拽平移/双击全屏预览图片 & 实时计算偏移量，长按450ms启动拖拽 \\
\cline{2-3}
& 进入分区，添加未分配物品 & 更新item.spaceId，关联物品到分区 \\
\cline{2-3}
& 替换图片/删除空间 & 确认对话框 → 级联清理分区和物品关联 \\
\hline
\multirow{4}{*}{事件管理} & 创建事件 + 勾选关联物品 & 写入event表 + event\_item关联表 \\
\cline{2-3}
& 查看事件详情与清单进度 & 查询event\_item，计算prepared/total \\
\cline{2-3}
& 点击清单项切换准备状态 & 更新isPrepared字段，实时刷新进度条 \\
\cline{2-3}
& 分享事件/标记完成 & systemShare生成格式化文本分享 \\
\hline
\end{tabular}
}
\end{center}

\blankline
\thirdheading{三、页面截图示意}

\begin{center}
\fbox{
\begin{tabular}{|>{\centering\arraybackslash}p{3.3cm}|>{\centering\arraybackslash}p{3.3cm}|>{\centering\arraybackslash}p{3.3cm}|>{\centering\arraybackslash}p{3.3cm}|}
\hline
\zihao{5}\bfseries 首页-物品列表 & \zihao{5}\bfseries 空间详情-标记点 & \zihao{5}\bfseries 事件详情-清单 & \zihao{5}\bfseries 个人中心-统计 \\
\hline
\multicolumn{4}{|c|}{\makecell[c]{
{\zihao{5}（顶部搜索栏 + 分类横向滚动标签 + 物品卡片列表）} \\
{\zihao{5}（空间平面图 + 彩色圆形标记点 + 分区列表）} \\
{\zihao{5}（标题/日期/进度条 + 物品清单（✅/⬜可切换）+ 分享/完成按钮）} \\
{\zihao{5}（物品总数 + 分类分布 + 事件统计 + 进行中提醒 + 背景设置入口）}
}} \\
\hline
\end{tabular}
}
\end{center}

\blankline

\secondheading{介绍文档}

\blankline

\thirdheading{一、创意背景}

据艾瑞咨询，智能家居行业用户规模达\textbf{1.9亿}，但应用场景单一、粘性持续下滑，亟需新场景。小红书社交平台调研显示，\textbf{72\%}的用户有物品数字化管理需求，\textbf{55\%}关注物品与空间的关系，\textbf{52\%}的收纳行为由生活事件驱动。结合问卷数据，\textbf{93.33\%}依赖记忆管理物品，\textbf{80\%}曾忘记物品位置。核心痛点在于缺乏直观的物品定位手段，且物品管理与事件准备之间缺少关联。"枣寻（LocateS）"以空间平面图可视化标记为核心、事件清单为驱动，填补智能家居场景空白，打造所见即所得的智能收纳管家。

\thirdheading{二、核心功能设计与技术实现}

\textbf{1. 物品管理：}支持多分类管理和名称搜索，可自定义分类。物品可分配至具体空间位置，显示层级化标签。\textbf{——实现：}category表（UNIQUE约束）存储分类，item表通过ItemRepository提供CRUD，AppStorage全局状态驱动跨页面自动刷新。

\textbf{2. 空间可视化标记：}用户上传平面图，点击图片放置彩色标记点创建分区，物品关联至分区，支持拖拽平移和全屏预览。\textbf{——实现：}space表以\verb|markerX|/\verb|markerY|归一化坐标（0–1）记录标记位置。clampImageOffset确保拖拽不越界，替换图片时级联清理子分区及物品关联。

\textbf{3. 事件驱动清单：}创建事件时从物品库勾选所需物品，生成准备清单。清单项可切换勾选状态，已完成项显示删除线，进度条实时更新，支持一键分享。\textbf{——实现：}event表与event\_item关联表实现多对多关系，EventRepository提供完整CRUD。采用Optimistic UI策略先更新本地再异步同步，分享集成\verb|systemShare|。

\textbf{4. 全局沉浸体验：}基于HDS组件库，底部悬浮页签配合沉浸光感材质，支持MiniBar展开/收起，提供6种渐变色和4种图片背景预设。\textbf{——实现：}HdsTabs配合\verb|setWindowLayoutFullScreen(true)|全屏。BackgroundStore持久化背景配置，PageBackground统一渲染。集成2×2/2×4服务卡片。

\thirdheading{三、市场前景}

智能家居赛道用户近\textbf{2亿}，但需求正从基础控制向更丰富场景延伸。小红书数据表明，物品记录（\textbf{72\%}）、整理收纳（\textbf{64\%}）、空间管理（\textbf{55\%}）需求旺盛，而市场上仍缺少将空间可视化、物品管理与事件驱动结合的智能收纳产品，处于蓝海阶段。\textbf{83.34\%}的问卷受访者对智能收纳持积极态度，学生宿舍群体占63.33\%。随着AI渗透率突破\textbf{54.7\%}，AI辅助整理、拍照分类、多设备协同等扩展功能可期，商业前景广阔。

\vspace{1.0cm}

\thirdheading{附：用户调研关键数据（n=30）}

\begin{center}
\begin{tabular}{|>{\centering\arraybackslash}p{5.2cm}|c|c|}
\hline
\zihao{5}\bfseries 调研维度 & \zihao{5}\bfseries 关键数据 & \zihao{5}\bfseries 对应功能 \\
\hline
依赖记忆管理物品 & 93.33\% & 物品管理系统 \\
\hline
收拾后忘记物品位置 & 80.00\% & 空间可视化标记 \\
\hline
每周至少找一次东西 & 46.67\% & 分类搜索+空间定位 \\
\hline
临近事件才想起准备 & 86.67\% & 事件驱动清单 \\
\hline
希望记录物品存放位置 & 76.67\% & 空间影像标记 \\
\hline
希望事件主动提醒 & 70.00\% & 事件清单提醒 \\
\hline
希望生成数字化地图 & 46.67\% & 空间图片+标记点 \\
\hline
\end{tabular}
\end{center}


% \clearpage

% \firstheading{二、格式及其他要求}

% （一）字体：宋体

% （二）字号

% 1、标题：二号，粗体

% 2、作者：三号，粗体

% 3、一级标题：三号，粗体

% \hspace{2em}二级标题：四号，粗体

% \hspace{2em}三级标题：小四号，粗体

% \hspace{2em}四级标题：五号，粗体

% 4、正文：五号

% （三）行距：单倍行距

% （四）页面设置

% 1、页边距：上：2.5厘米\hspace{1em}下：2.5厘米

% \hspace{4em}左：3厘米\hspace{2em}右：3厘米\hspace{1em}装订线：0厘米

% 2、页眉：1.5厘米

% \hspace{2em}页脚：1.5厘米

% 3、纸型：A4，纵向

% （五）插入页码

% \hspace{2em}位置：页面底端

% \hspace{2em}对齐方式：外侧

% （六）其他

% 1、PDF格式提交，命名格式为 01-作品说明文档+参赛队伍名称

% 2、文档中若包含插图、表格请按序编号并命名

% 3、主体内容不超过20页（参考资料及附录不计入总页数）

\end{document}
